\documentclass[11pt]{article}
\usepackage[a4paper,margin=1in]{geometry}
\usepackage{amsmath,amssymb,mathtools,bm}
\usepackage{enumitem,booktabs,array}
\usepackage{tcolorbox}
\usepackage{hyperref}
\usepackage[protrusion=true,expansion=false]{microtype}
\hypersetup{colorlinks=true,linkcolor=blue,urlcolor=blue,citecolor=blue}
\setlist[itemize]{topsep=2pt,itemsep=2pt}
\setlist[enumerate]{topsep=2pt,itemsep=2pt}
\newtcolorbox{keybox}{colback=blue!3!white,colframe=blue!40!black,boxrule=0.4pt,arc=1mm}
\newtcolorbox{alertbox}{colback=red!3!white,colframe=red!40!black,boxrule=0.4pt,arc=1mm}
\newtcolorbox{answerbox}{colback=green!3!white,colframe=green!40!black,boxrule=0.4pt,arc=1mm}
\newtcolorbox{drillbox}{colback=yellow!5!white,colframe=orange!60!black,boxrule=0.4pt,arc=1mm}
\newcommand{\EV}{\mathbb{E}}
\newcommand{\Var}{\mathrm{Var}}
\newcommand{\Cov}{\mathrm{Cov}}
\newcommand{\blank}[1]{\underline{\hspace{#1}}}

\title{E02-01: Brav, Jiang \& Kim (2015, RFS) \\
\large ``The Real Effects of Hedge Fund Activism'' --- Active Recall Workbook}
\author{Empirical Corporate Finance II --- Exam Prep}
\date{\today}
\begin{document}\maketitle

\begin{keybox}
\textbf{Paper in one sentence.} Using US Census plant-level data matched to hedge-fund 13D activism events 1994--2007, BJK show that target firms' plants see productivity gains ($\sim$5--6\% within 3 years), labour-intensity declines, and operating margins rise --- direct evidence that activism produces \emph{real} operational improvements, not just financial engineering.

\textbf{Placement.} Canonical empirical paper on hedge-fund activism. Paired with Albuquerque-Fos-Schroth (E02-02) which structurally models value creation in activism. Contrast with short-termism critiques.
\end{keybox}

\section*{Round 1: Complete Walk-Through}

\subsection*{1.1 Data and setting}
\begin{itemize}
\item US Census Longitudinal Business Database and Census of Manufactures plant data (establishment level).
\item Hedge-fund 13D filings 1994--2007 identifying activist targets ($\sim$2{,}000 events).
\item Match target firms to plant-level outcomes.
\end{itemize}

\subsection*{1.2 Empirical design}
Plant-level DiD:
\[
Y_{pt}=\alpha_p+\delta_t+\sum_{\tau=-3}^{+3}\beta_\tau\,\mathbf 1[t-\text{event}=\tau]+X_{pt}'\gamma+\varepsilon_{pt}.
\]
Outcomes: TFP, labour productivity, capital intensity, wages, employment, margins. Control: matched plants from non-targeted firms.

\subsection*{1.3 Headline results}
\begin{itemize}
\item Plant-level TFP rises $\sim$5--6\% over 3 years post-activism.
\item Labour productivity rises more; capital intensity falls --- substitution away from labour.
\item Wages per worker rise modestly; employment falls; total payroll unchanged.
\item Plants closed post-activism are below-average in TFP; survivors are stronger.
\item Effects concentrated in ``offensive'' campaigns (strategic/M\&A) vs.\ ``financial'' campaigns.
\end{itemize}

\subsection*{1.4 Mechanism}
\begin{itemize}
\item Activists push targets to reallocate resources toward more-productive uses.
\item Closing unproductive plants; improving remaining ones.
\item Monitoring reduces managerial slack, especially in diversified firms.
\end{itemize}

\subsection*{1.5 Interpretation}
\begin{itemize}
\item Activism creates real value, not just price appreciation via expropriation.
\item Short-termism critique (that activists extract wealth at expense of long-run investment) is not supported at the plant level.
\item Labour losses are real but are concentrated in uncompetitive plants.
\end{itemize}

\subsection*{1.6 Caveats}
\begin{alertbox}
\begin{itemize}
\item 13D filings are selective: activists target firms they expect to improve.
\item Plant-level measurement of TFP has well-known issues (mark-ups, outsourcing).
\item No effects on R\&D / innovation pipelines beyond 3 years examined.
\end{itemize}
\end{alertbox}

\section*{Round 2: Level 1 Fill-in}
\begin{enumerate}
\item Data: US \blank{1cm} plant data + hedge-fund \blank{1cm}D filings 1994--\blank{1cm}.
\item Plant-level TFP rises $\sim$\blank{1cm}\% within 3 years.
\item Employment \blank{1cm}s; wages per worker \blank{1cm}.
\item Closed plants: \blank{1cm} than firm average in TFP.
\item Campaigns with biggest effects: \blank{1.5cm}.
\end{enumerate}

\section*{Round 3: Level 2 Prompts}
\begin{enumerate}
\item Describe the plant-level DiD and identification concerns.
\item Discuss the TFP, labour, and employment results.
\item How does BJK refute the short-termism critique of activism?
\item Relate to Albuquerque-Fos-Schroth (2022).
\item What are the identification threats, and how does the paper address them?
\end{enumerate}

\section*{Round 4: Minimal Scaffolding}
From a blank page: (i) data + design, (ii) TFP results, (iii) labour results, (iv) mechanism, (v) implications.

\section*{Round 5: Blank Page}
Essay: does hedge-fund activism create real value? BJK's plant-level evidence.

\vspace{6pt}
\begin{drillbox}
\textbf{Speed Drill.} (1) Data sources. (2) TFP effect size. (3) Employment effect. (4) Closed-plant selection. (5) Interpretation vs.\ short-termism.
\end{drillbox}

\section*{Quick Reference \& Answer Key}
\begin{answerbox}
\begin{itemize}
\item \textbf{Data.} US Census plants + 13D filings, 1994--2007.
\item \textbf{TFP.} $+5$--$6$\% within 3 years.
\item \textbf{Labour.} Employment falls; wages rise.
\item \textbf{Lesson.} Activism generates real operational gains.
\end{itemize}
\end{answerbox}

\begin{keybox}
\textbf{30-second exam pitch.} Brav, Jiang \& Kim (2015) use US Census plant-level data matched to hedge-fund 13D activism events (1994--2007, $\sim$2{,}000 events) to test whether activism creates real operational value. Plant-level DiD shows treated plants gain 5--6\% TFP within three years, labour productivity rises, capital intensity falls, wages per worker rise modestly, employment falls, and plants closed post-activism are below-average in productivity --- consistent with reallocation toward more-productive uses. ``Offensive'' campaigns (strategic/M\&A) drive the largest effects. The results counter the short-termism critique that activists extract wealth at the expense of long-run productivity; the real effects are improvements, not wealth transfers. BJK is the canonical empirical demonstration that activism has real consequences in production, and it frames the 2010s-20s debate on activist investing and corporate governance.
\end{keybox}

\end{document}
